\documentclass[../main.tex]{subfiles}
\begin{document}
%==========================================TIÊU ĐỀ TẬP========================================
\begin{table}[h]
  \begin{adjustwidth}{-1cm}{}
    \begin{tabular}{>{\centering\arraybackslash}p{6cm}|>{\raggedright\arraybackslash}p{12.5cm}}
      \multirow{5}{6cm}
      % Cột bên trái: ÉPISODE và số tập
      {\\[-40pt]
        \flaregothic\fontsize{20pt}{20pt}\selectfont \centering
        \color{eptype}ÉPISODE\color{black}\\
        \burbank\fontsize{72pt}{72pt}\selectfont
        \color{epnum}59\color{black}\\[6pt]
        \cafeteria\fontsize{20pt}{20pt}\selectfont\color{epnum}(5-3)\color{black}
        \\[18pt]
        \color{black}
      }
       &
      % Dòng thứ nhất: tên tập tiếng Nhật
      {\fotskip\fontsize{18pt}{18pt}\selectfont \ruby{全}{ぜん}\ruby{力}{りょく}で\ruby{追}{お}いかけろ！
      } \\[6pt]
      % Dòng thứ hai: tên tập tiếng Anh
       & {\cafeteria\fontsize{18pt}{18pt}\selectfont Lamborg*ini} \\[4pt]
      % Dòng thứ ba: tên tập tiếng Việt
       & {\vnmsans\fontsize{18pt}{18pt}\selectfont Bài toán quãng đường} \\[8pt]
      % Dòng thứ tư: Nhân vật xuất hiện
       & {\caviar{\fontsize{14pt}{14pt}\selectfont Nhân vật: \emp{kuon1} \emp{aki} \emp{fubuki} \emp{matsuri} \emp{choco} \emp{subaru} \emp{mio} \emp{okayu} \emp{korone}}}
      \\[8pt]
      % Dòng cuối cùng: Ngày phát hành
       & {\continuum\fontsize{16pt}{16pt}\selectfont Ngày phát hành: 21/06/2020} \\[18pt]
    \end{tabular}
  \end{adjustwidth}
\end{table}
%=========================================TRUYỆN CHÍNH========================================
\color{black}

% Hộp tóm tắt
\txtbx[2]{{\color{vol} \uline{Tóm tắt:}} Matsuri vô tư rời văn phòng mà quên hết những món đồ quan trọng. Fubuki hoảng hốt, Mio buộc phải chạy đuổi giao đồ. Trên đường, Kuon nhập cuộc, rồi Okayu xuất hiện với chiếc ``Lamb*rgh*ni tên lửa'' chở theo Korone, Subaru và Aki. Một ``vị thần'' mang hình hài Choco ra bài toán vận tốc Mach 80-100 khiến cả nhóm quay cuồng. Kết thúc, trên xe hỗn loạn vì cướp đồ nghề của ``thần''.}

\begin{center}
  \uline{(Người viết tập này có sử dụng bản dịch từ \href{https://www.youtube.com/@TomangclipHololivevedich}{Tớ mang clip \hll\ về dịch - TMCHVD}.)}
\end{center}

``Tôi đi nha!''

Matsuri ló mặt ra khỏi cửa, tay vẫy vẫy như một đứa trẻ sắp đi chơi xa. Nụ cười trên môi cô rạng rỡ, chẳng chút muộn phiền.

``Bảo trọng!'' Fubuki và Mio vui vẻ đáp lại, giọng đầy ấm áp.

\begin{figure}[H]
  \centering
  \begin{subfigure}[t]{0.45\textwidth}
    \centering
    \includegraphics[width=8cm]{../../../../Image/Book5/5-3a1.png}
  \end{subfigure}
  \quad
  \begin{subfigure}[t]{0.45\textwidth}
    \centering
    \includegraphics[width=8cm]{../../../../Image/Book5/5-3a2.png}
  \end{subfigure}
\end{figure}

Cánh cửa đóng sầm lại. Và rồi, bầu không khí yên bình trở lại, hai cô nàng thuộc nhóm GAMERS - Fubuki và Mio - ở lại văn phòng tiếp tục công việc của mình. Tiếng bàn phím lách cách, tiếng giấy lật sột soạt, tất cả tạo nên một khung cảnh thật thư thái.

Về phía Matsuri, hôm nay cô có việc đột xuất nên phải ra ngoài trước. Còn Kuon thì vừa lúc ra ngoài mua chút đồ dự trữ cho văn phòng. Mọi thứ có vẻ diễn ra theo đúng kế hoạch.

Thế nhưng, chỉ mới đi được một lúc, nàng Cáo trắng bỗng giật mình. Đôi mắt cô mở to, nhìn chằm chằm vào đống đồ đạc vụn vặt nằm ngổn ngang trên bàn - chỗ Matsuri vừa ngồi cách đây chưa đầy một phút.

``Chết rồi!'' cô thốt lên, giọng đầy bàng hoàng - ``Nó để quên chìa khóa nhà, điện thoại, ví và cả thẻ PASMO\footnote{Một loại thẻ điện tử trả trước dùng khi đi tàu điện ngầm.} rồi!''

\img{5-3b}

Nàng Sói đen nhìn đống đồ đạc kia, lắc đầu ngán ngẩm. Cô chống cằm, giọng đầy bất lực: ``Chắc liệt kê mấy cái chị ấy có mang đi còn nhanh hơn đấy\dots\ Ví, điện thoại, chìa khóa - ba thứ quan trọng nhất của một con người hiện đại, thế mà chị ấy bỏ hết lại. Chị ấy có mang theo thứ gì không nhỉ?''

Về cơ bản, Matsuri đã rời khỏi văn phòng mà không mang theo bất cứ thứ gì cần thiết. Cô nàng ra đi trong sự phấn khích, để lại sau lưng một đống tài sản cá nhân đang nằm chờ bị bỏ rơi.

Thấy tình hình có vẻ không ổn, hai người liền bàn bạc tìm cách giải quyết.

Fubuki - với vẻ mặt đầy lo lắng - lên tiếng trước: ``B-Bọn mình cần phải nhanh chóng mang đồ đến cho nó thôi! Không có điện thoại, bà ấy không thể gọi ai; không có ví, Matsuri không thể mua gì; không có chìa khóa, tối nay bà ấy ngủ ngoài đường mất!''

Mio bình tĩnh đáp lại, giọng lý trí: ``Nhưng mà này, nếu chị ấy để quên nhiều như vậy thì kiểu gì lát nữa cũng sẽ tự quay lại lấy thôi mà? Con người ta đi được một đoạn, sờ túi thấy thiếu đồ là quay về ngay.''

Fubuki lắc đầu quầy quậy, mắt nhìn xa xăm: ``Không đâu! Nếu là bà ấy thì\dots'' cô vừa nói vừa tưởng tượng ra cảnh tượng trong đầu.

Trong tâm trí nàng Cáo trắng, Matsuri đang đứng chết trận giữa đường, mặt đờ đẫn như một con cá ươn vừa bị kéo lên khỏi mặt nước. Cô nàng đưa tay sờ túi, rồi sờ túi kia, rồi lục tung khắp người, rồi đứng như trời trồng khi nhận ra mình chẳng có gì. Rồi cô tự trách móc cái đầu óc cá vàng của mình, vò đầu bứt tai, có khi còn đấm vào tường vì xấu hổ.

``Chắc giờ nó đang kêu cứu thảm thiết ở đâu đó rồi!'' Fubuki kết luận, giọng như sắp khóc.

\img{5-3c}

Thậm chí, nàng Cáo trắng còn có thể ``nghe thấy'' tiếng kêu cứu của Matsuri vọng ra từ đằng xa, mặc dù khoảng cách thực tế chắc phải cỡ\dots\ cả cây số. Cô nghiêng đầu, ra hiệu cho Mio lắng nghe.

Mio ngơ ngác, chẳng nghe thấy gì ngoài tiếng xe cộ ngoài phố: ``Nhưng mà này, sao bà nghĩ chị ấy sẽ không quay lại đây? Chị ấy mới chỉ chạy đi chưa đầy một phút thôi mà. Quay lại cũng chỉ mất hai phút chứ nhiêu.''

Fubuki đáp lại đầy quả quyết, mắt long lanh như thể vừa thực hiện một phép tính siêu việt: ``Bởi vì bà ấy đã đi xa lắm rồi!''

Cô nheo mắt, đưa tay lên tai như một radar sống, lắng nghe một lần nữa. Rồi cô chỉ tay thẳng về hướng Nam, giọng dõng dạc: ``Khoảng\dots\ 500 mét theo hướng đó! Đã qua hai ngã tư, sắp đến cửa hàng tiện lợi. Tốc độ của nó là khoảng\dots\ 7 km/h - đi bộ nhanh lắm đấy!''

Mio nhìn cô bạn, mắt chớp chớp: ``\dots\ Bà đang đùa đấy chắc? Làm sao bà biết chính xác được?''

Fubuki nghiêm túc đáp lại, giọng đầy tự tin đến mức đáng ngờ: ``Giờ phải có ai đó đuổi theo rồi trả đồ lại cho nó ngay! Chậm một giây là thêm một trăm mét, chậm một phút là lạc luôn dấu vết!''

Nàng Sói đen thở dài, biết rằng chẳng thể cãi lại được cơn ``hỏa tốc'' của nàng Cáo. Cô nhìn đống đồ của Matsuri, rồi nhìn ra cửa, rồi nhìn Fubuki: ``Được rồi, vậy ai sẽ là người chạy đi đưa đây?''

\ct

``\textbf{Tại sao lại phải là mình chứ!?}''

\gif{5-3d}{1}{66}

Mio chạy như điên, lòng đầy ức chế, miệng không ngừng càu nhàu. Mồ hôi bắt đầu lấm tấm trên trán, nhưng cô chẳng dám dừng lại. Trong tay nàng Sói đen là đống đồ của Matsuri, tất cả những thứ mà một con người bình thường không thể ra khỏi nhà nếu thiếu.

``Sao cái số mình lại khổ thế này? Đang yên đang lành ngồi trong văn phòng, tự dưng phải đi giao hàng tốc hành cho cái người mà suốt ngày chỉ biết la hét `Wasshoi! Wasshoi!' chứ\dots''

Câu nói của nàng Cáo trắng Fubuki lúc này cứ văng vẳng trong đầu cô lúc chia tay: ``Cảm ơn bà nhé, Mio\~{}'' ngọt ngào đến mức làm cô càng thêm bực mình. ``Cám ơn gì mà cám ơn! Sao bà không tự đi mà đưa?''

Bất chợt, một giọng nói cất lên ngay bên cạnh, điềm nhiên như đang đi dạo trong công viên: ``Thì bởi vì chính bản thân Sói chạy khá nhanh mà.''

Mio giật mình quay sang. Kuon bất thình lình xuất hiện, chạy song song với cô, hơi thở vẫn đều đều, mặt không chút mệt mỏi, như thể đang chạy bộ buổi sáng chứ không phải đuổi theo một kẻ chạy trốn.

Nàng Sói đen trừng mắt nhìn cậu quản lý: ``Em không có nói thế! Mà anh tự dưng ở đâu xuất hiện vậy!? Có phải anh lén theo dõi em không?''

Cậu nhún vai, giọng tỉnh bơ: ``Đang trên đường từ cửa hàng tiện lợi về phòng thì thấy em chạy như ma đuổi, tưởng có chuyện gì nên chạy theo hỏi đây. Có vụ gì mà long trời lở đất thế?''

Mio tiếp tục lao về phía trước, mặt nhăn nhó như mếu: ``Matsuri-senpai quên đồ, nên em mới chạy đưa lại cho chị ấy! Chị ấy đi rồi, chẳng mang theo cái gì cả. Mà chị ấy cũng nhanh thật, mới có vài phút mà đã xa tít rồi.''

Kuon khẽ ngạc nhiên: ``\textit{Matsuri mà quên đồ ư? Có gì lạ đâu, cô nàng vốn nổi tiếng là đầu óc trên mây. Nhưng quên nhiều đến thế này thì cũng hơi\dots\ ấn tượng.}'' Cậu hỏi, giọng đầy vẻ hoài nghi: ``Quên gì?''

Mio thở dài, liệt kê tăm tắp những thứ cô đang cầm trên tay như một người bán hàng rong: ``Ối giồi! Điện thoại, ví, thẻ PASMO, chìa khóa nhà\dots\ Nói chung là nhiều lắm! Tóm lại là tất cả những thứ cần thiết cho sự sống còn của một con người hiện đại - chị ấy để hết lại bàn. Chỉ thiếu mỗi cái quần áo trên người thôi.''

Cậu con trai nghe vậy, lắc đầu ngán ngẩm. Trong đầu cậu, một ý nghĩ vụt qua: ``\textit{Thật không thể hiểu nổi, người như Matsuri mà lại có thể tồn tại đến ngày hôm nay\dots}''

Cậu bật ra câu châm chọc, giọng đầy vẻ ``thương hại'': ``Em ấy óc nho đến thế à? Thế sao không quên luôn cả thân mình đi cho rồi? Để khỏi phải chạy đi chạy lại.''

Mio cảm thấy máu trong người dồn lên não. Cô gần như hét lên, tức giận không thể kiềm chế: ``\textbf{Anh còn đùa được thì em cũng chịu! Mau cầm hộ em đồ giùm cái! Tay em sắp rụng rồi!}''

Cô vừa nói vừa ném đống đồ của Matsuri lên cao, nhắm thẳng về phía cậu Tổng đang chạy ngay phía sau. Kuon phản xạ nhanh, giơ tay bắt gọn lấy đống đồ như một thủ môn bắt bóng, không cần giảm tốc độ. Hai người tiếp tục lao đi như hai vận động viên điền kinh trong một cuộc thi không chính thức.

Đột nhiên, một luồng sáng kỳ lạ lóe lên bên cạnh họ. Nó không gay gắt, mà nhẹ nhàng, huyền ảo, như thể ai đó vừa mở một cánh cửa dẫn từ thế giới khác sang.

Từ trong ánh sáng đó, một bóng dáng quen thuộc hiện ra. Đó là một người phụ nữ với dáng vẻ điềm nhiên, khoan thai, mái tóc dài bay nhẹ trong gió. Trông cô ấy giống hệt Yuzuki Choco - nàng Y tá ngọt ngào của \hll - nhưng có gì đó khác. Một cây trượng dài trên tay, một đám mây nhỏ dưới chân nâng cô lơ lửng trên không. Cô mỉm cười bí ẩn, rồi cất tiếng, giọng vừa quen vừa lạ: ``Ta có bài cho hai con giải đây.''

\img{5-3e}

Cả hai đang mải chạy bỗng giật mình, suýt vấp ngã. Mio thảng thốt: ``C-Choco-sensei!?''

Kuon cũng ngỡ ngàng: ``Choco-san? Sao em lại cưỡi trên mây thế kia? Mới sáng ra đã uống thuốc gì thế?''

Người phụ nữ lắc đầu, nụ cười vẫn giữ nguyên, nhưng giọng có chút trêu chọc: ``Không phải đâu à nghen. Nhầm rồi. Ta không phải Choco mà các con biết. Ta là một vị thần lang bạt trên thế gian, muốn thử thách những con người có trí tuệ nhanh nhẹn nhất.''

Mio nhíu mày, vừa chạy vừa liếc sang với ánh mắt đầy nghi hoặc: ``Người không phải là cô ấy sao? Giống hệt từ A đến Z, từ giọng nói đến kiểu cười. Có khi nào cô đang cosplay thần không?''

Kuon thì lại khác. Cậu tặc lưỡi, phán luôn một câu đầy thẳng thắn, chân vẫn không ngừng bước: ``Anh chắc đó là em ấy đấy. Chẳng qua đang rảnh quá nên nghĩ ra trò. Mà, bài gì, ra đề đi! Đang chạy gấp lắm, không có thời gian chờ đợi đâu!''

Nàng Sói đen chưa kịp phản ứng trước sự chấp nhận nhanh chóng của cậu con trai thì ``vị thần'' y tá kia đã bắt đầu ra đề, giọng đều đều như giáo viên đang đọc bài giảng: ``Gọi P1, tức Matsuri-sama, rời đi vào lúc 10 phút trước. Còn hai con, gọi là P2, đang đuổi theo để trả lại mấy món đồ mà người đó để quên.''

\img{5-3f}

Cậu Tổng nhướng mày, tiếp lời: ``Vậy coi vận tốc của bọn con tính theo người bình thường - khoảng 10-12 km/h, chứ gì? Chạy nhanh hơn nữa thì chết sớm.''

``Ta cho là vậy. Giờ hãy tính thời gian hai con cần để đuổi kịp Matsuri-sama.''

Mio nhăn mặt, thầm nghĩ: ``\textit{Ngay trong lúc khẩn cấp này mà cô vẫn phải làm toán sao!? Chúng ta đang giải cứu một kẻ mất trí chứ không phải đi thi đại học!}''

Cô kêu lên đầy bất bình, giọng vang vọng cả con phố: ``Bộ đây là chương trình giáo dục trực tuyến à!? Giờ thể dục sao lại có toán ở giữa thế?''

Cậu Tổng thì điềm nhiên đáp, chân vẫn không ngừng chạy, mắt vẫn nhìn về phía trước: ``Thì cứ coi là như vậy đi. Thế Mat-chan chạy với tốc độ bao nhiêu? Đưa số liệu cụ thể vào, không thì giải bằng ẩn à?''

Nhưng ngay lúc đó, ``vị thần'' kia tiếp tục, giọng vẫn điềm đạm, như thể đang kể một câu chuyện cổ tích: ``Biết rằng Matsuri-sama di chuyển với vận tốc Mach 80.''

Cả Mio và Kuon lập tức khựng lại.

Không phải là chậm dần, mà là đứng sững như trời trồng. Mắt họ mở to, miệng há hốc. Một khoảng lặng ngắn ngủi, rồi cả hai đồng thanh hét lên, giọng xé toạc bầu không khí buổi sáng: ``\textbf{Sao mà nhanh vậy!?}''

Nàng Sói đen há hốc mồm, quên cả chạy: ``Tưởng chị ấy chỉ đi bộ thôi sao!? Mach 80 là gấp 80 lần tốc độ âm thanh - gần 100.000 km/h! Cả tên lửa còn chưa nhanh bằng!''

Cậu Tổng cũng đổ mồ hôi hột, tính nhẩm trong đầu: ``Với vận tốc đó, em ấy đã bay qua Thái Bình Dương mấy lần rồi! Không, em ấy có thể vòng quanh Trái Đất trong vòng\dots\ hai mươi phút? Mà sao chúng ta chưa thấy bóng hồng nào vụt qua chứ?''

``Vị thần'' Y tá (hay Choco, hay ai đó) vẫn giữ nụ cười nhẹ nhàng, trả lời như thể đó là chuyện hiển nhiên, ai cũng phải biết: ``Không, cô ấy lái Lamborg(h)ini đấy. Xe thể thao, tốc độ tối đa có thể lên đến\dots\ ờm, trên lý thuyết là Mach 80 nếu có cánh và đốt nhiên liệu tên lửa.''

Đến đây, cậu Tổng trợn mắt, hét lên ngay lập tức, giọng đầy bất lực: ``\textbf{Thế thì bố ai mà đuổi kịp nổi!? Gọi không quân đi! Hay gọi Superm*n đi!}''

``Vị thần'' Y tá chỉ nhún vai, đáp lại một câu ngắn gọn nhưng đầy tính triết lý, như thể đang khai thị cho hai kẻ trần mắt: ``Đúng rồi đấy\~{}''

Mio đứng chết trân, hai tay buông thõng, đống đồ của Matsuri trên tay Kuon bỗng trở nên vô nghĩa. Cô cảm thấy mình sắp phát điên. Tình huống này chẳng khác gì một trò đùa của thượng đế: người ta chạy bằng tốc độ âm thanh, còn mình chạy bằng chân - đuổi sao lại?

Cô hét lên đầy tuyệt vọng, ngửa mặt lên trời: ``Làm ơn quay lại cho bọn em nhờ, Matsuri-senpai! Bọn em bắt đầu phát điên lên rồi đấy!''

\ct

Sau câu nói đầy chán nản của cậu con trai, cả hai gần như mất hết động lực chạy tiếp. Nàng Sói đen tuyệt vọng, định ngửa mặt lên trời mà hét lên một lần nữa cho hả dạ, thì bỗng một giọng nói vang lên từ phía xa, lạnh lùng nhưng đầy\dots\ cứu rỗi: ``Cần tôi giúp không?''

Cả hai lập tức quay đầu lại. Và rồi, họ chứng kiến một cảnh tượng mà có lẽ cả đời cũng khó quên.

Không ai khác, chính là Okayu, người đang lái một chiếc xe\dots\ cực kỳ quái dị. Nó có tạo hình giống một quả tên lửa siêu thanh màu xám, thon dài, khí động học, nhưng lại đủ chỗ cho năm người ngồi. Và quan trọng nhất: nó không có bánh xe.

\img{5-3g}

\dots\ Thực ra, nó chính là quả tên lửa đã được cải biên từ vụ đua Tết năm ngoái. Đây cũng chính là chiếc xe mà nàng Mèo đã dùng để đua với Fubuki trong dịp Tết vừa rồi\footnote{Xem lại tập đặc biệt Tết.}, chỉ khác là lần này nó được trang bị thêm ghế ngồi và một cái vô lăng giả.

Mio trợn mắt trước sự xuất hiện của chiếc xe-không-ra-xe này, giọng run run: ``Cái\dots\ cái gì đây!? Tên lửa à? Máy bay không cánh? Hay là món đồ chơi của một tỷ phú nào đó?''

Okayu đáp lại bằng vẻ mặt vô cùng điềm tĩnh, giơ ngón tay hình chữ V, giọng đầy tự hào: ``Lamborg*ini đó!''

Ngay lập tức, nàng Sói đen gắt lên đầy bực tức, giọng như xé gió: ``\textbf{Đừng có xạo!} Lamb*rgh*ni mà không có bánh, lại còn bay lơ lửng thế này à!? Đây là Lamb*rgh*ni hay là UFO?'' Rõ ràng, Mio đã quá chán ngán khi cứ liên tục nghe cái tên hãng xe đến từ nước Ý này trong suốt mười lăm phút qua.

Kuon thì có vẻ tò mò hơn. Cậu trèo lên, nhìn chiếc xe một lượt từ đầu đến cuối, rồi hỏi bằng giọng của một kỹ sư hàng không: ``Sao trông nó giống hệt tên lửa siêu thanh thế? Có cánh, có đuôi, có cả ống xả phản lực đây này. Em chế từ quả tên lửa thật à?''

Okayu nhún vai, thản nhiên nói tiếp, như thể đang giới thiệu một mẫu xe mới ra mắt thị trường: ``Nó có thể phóng với vận tốc Mach 100 luôn đó. Đi từ Tokyo đến New York mất chưa đầy ba phút.''

Mio gần như nghẹn lời, hai tay bóp đầu: ``\textbf{Lại Mach à!!??} Tôi chán cái từ này rồi! Mach, Mach, Mach - cái gì cũng Mach! Có khi tôi cũng biến thành Mach mất thôi!''

Cậu con trai ở bên cạnh thì chỉ có thể cố gắng an ủi nàng Sói đang dần mất kiên nhẫn, vỗ vai cô: ``Thôi, bình tĩnh. Nhanh hơn Mach 80 là được rồi. Miễn đuổi kịp Mat-chan là được.''

Sói đen không còn muốn tranh luận thêm nữa. Cô phất tay đầy bất lực, như một chiến binh đầu hàng trước số phận: ``Kệ đi! Cứ cho chúng tôi đi nhờ là được! Có xe gì thì đi xe đó, có tên lửa thì ngồi tên lửa - miễn sao giao được đồ cho cái người đãng trí kia là được!''

\ct

Bên trong chiếc ``tên lửa siêu thanh'', không gian chật chội hơn Mio tưởng rất nhiều. Cô ngồi ở ghế trước, cạnh Okayu đang ung dung cầm vô lăng. Kuon thì bị đẩy ra hàng ghế sau cùng, ngồi co ro giữa hai dãy ghế, trông như một hành khách hạng ba trên tàu hỏa thời xưa.

Ở hàng ghế giữa, ba con người đang ngồi một cách điềm tĩnh, mặt thản nhiên như đang đi xe buýt: Korone đang nhai kẹo, mắt nhìn ra cửa sổ; Subaru tay ôm balo, mặt hơi xanh vì tốc độ; Aki đang lấy điện thoại ra selfie, vì ``lần đầu được ngồi tên lửa''.

Okayu vẫn ung dung giữ vô lăng, điều khiển con ``tên lửa siêu thanh'' lướt trên đường với một vẻ mặt đầy thư thái, như thể đang lái xe đạp dạo quanh công viên.

\img{5-3h}

Mio thở hổn hển một lúc, lấy lại hơi thở, rồi bắt đầu cảm thấy khó chịu. Cô nhăn mặt, vặn vẹo người, than phiền: ``T-Tự nhiên trong này chật quá vậy\dots\ Năm người ngồi trên một cái tên lửa vốn chỉ có hai chỗ - sao mà thở nổi! Chân em tê hết rồi!''

Từ phía sau, cậu Tổng lên tiếng, giọng điềm nhiên nhưng đầy chua chát: ``Có xe đi nhờ là tốt lắm rồi. Đừng đòi hỏi làm gì. Đời không như là mơ.''

Ngay lập tức, nàng Sói quay ngoắt lại, mắt trợn ngược, hét lớn như sấm: ``\textbf{Tại ai mà chúng ta rơi vào tình cảnh này hả!? Nếu không phải tại Matsuri-senpai quên đồ, thì chúng ta đã chẳng phải đi giao hàng tốc hành, mà cũng chẳng cần ngồi nhét mình trong cái thứ tên lửa nhà quê này!}''

Cậu Tổng đáp lại, giọng vẫn tỉnh bơ, nhưng có chút\dots\ đắc ý: ``Em ngồi ở hàng trước, xe còn có chỗ mà thở đấy. Anh còn bị đẩy ra đằng sau, bị mấy đứa ngồi ghế giữa chèn ép hết rồi đây này. Cái ghế này nó nhọn hoắt, tựa như ghế tra tấn thời trung cổ.''

Ba người ngồi hàng giữa lặng lẽ nhìn nhau, rồi chỉ còn nước đóng vai ``tổ hòa giải'', cố gắng dỗ dành hai người kia bớt cãi nhau.

Subaru lên tiếng trước, giọng đầy thiện chí: ``Thôi nào, hai người bình tĩnh. Quan trọng là chúng ta sắp đuổi kịp Matsuri-chan rồi.''

Korone nhai kẹo, gật đầu: ``Ừ ừ, cố lên, sắp đến nơi rồi, gâu!''

Aki vẫn cười toe: ``Em chụp được tấm selfie rất đẹp, lát nữa chia sẻ cho mọi người nhé\~{}''

Đúng lúc đó, Okayu đột nhiên chỉ tay về phía trước, mắt sáng lên, reo lên như một đứa trẻ vừa tìm thấy kho báu: ``Tìm thấy Matsuri-chan rồi!''

Lập tức, Mio quay ngoắt đầu về hướng nàng Mèo chỉ. Và quả thực, trên đường, cách đó không xa, Matsuri đang đứng như trời trồng, hai tay thõng xuống, ánh mắt vô hồn, mặt đờ đẫn như một con cá ươn vừa bị kéo lên khỏi mặt nước. Đúng y hệt lời tiên đoán của Fubuki ban nãy.

\img{5-3i}

Cảm giác vui mừng dâng trào. Mio như sắp khóc, cô quay sang Okayu, giọng đầy cảm kích: ``Tốt quá! Cảm ơn em! Cho bọn chị xuống--''

Cô chưa kịp nói hết câu thì Okayu bất ngờ cắt ngang, giọng thản nhiên như thể đang báo giá một món hàng: ``Xin lỗi, nhưng cái này không thể dừng gấp như vậy được.''

Khoảnh khắc đó, chiếc xe lao vút qua Matsuri như một cơn gió nhanh, mạnh, và không chút do dự. Cả bóng dáng cô nàng tóc nâu lùi dần phía sau, nhỏ xíu rồi biến mất.

\gif{5-3j}{1}{51}

Matsuri, vẫn trong trạng thái vô cảm, chỉ kịp lẩm bẩm một câu đầy hoang mang khi chiếc xe phóng ngang: ``\dots\ Cái gì vậy\dots?''

Bên trong xe, Mio hoảng loạn hét lên, hai tay bám vào thành ghế: ``\textbf{Quay lại! Quay lại ngay!! Chúng ta vừa đi ngang qua mất tiêu rồi!}''

Còn cậu Tổng thì lại rất ung dung, ngồi tựa vào ghế, bình thản buông một câu như đổ thêm dầu vào lửa: ``Anh tưởng vừa nãy em dừng xe được mà nhỉ?''

Nàng Mèo nhún vai, giọng vô cùng tự nhiên: ``Đấy là vì tôi đoán được Mio-chan sẽ ở đó thôi\~{} Khi có dữ liệu trước, tôi có thể tính toán quỹ đạo dừng. Còn lần này thì\dots\ bất ngờ quá.''

Mio càng điên hơn, mặt đỏ bừng: ``\textbf{Vậy mà sao giờ cô lại không thể dừng như vậy được chứ!? Cũng là bất ngờ, cũng là nhìn thấy, cũng là phanh - có gì khác!?}''

Okayu nhếch môi cười nhẹ, giọng đầy triết lý, như một giáo sư vật lý đang giảng bài: ``Không có số liệu để căn thời gian đâu. Khi gặp Matsuri-chan lần đầu, tôi đã có sẵn thông số. Còn bây giờ, xe đang ở tốc độ Mach 100, phanh gấp thì văng hết ra ngoài.''

Sói đen ngã quỵ xuống ghế, hai tay ôm mặt: ``Thế lại phải giải toán à!? Lại liên quan đến Mach, đến quãng đường, đến thời gian ư!? Đời tôi sao khổ thế này!''

Kuon gật đầu không chút do dự, giọng đầy vẻ thông cảm: ``Chứ sao nữa em. Liên quan đến vật lý cả đấy. Nhưng mà nếu Okayu-chan không tính được, thì chẳng ai tính nổi đâu.''

Nàng Sói lại gào lên, giọng xé toạc cả khoang xe: ``\textbf{Tôi không cần biết! Tôi không cần toán, không cần vật lý, không cần Mach! Mau quay lại đi!!}''

Nàng Mèo vẫn giữ vẻ mặt thư thái, chỉ nhún vai rồi phán một câu khiến Sói đen sụp đổ hoàn toàn: ``Cái đó thì tôi chịu ha\~{} Vật lý là vật lý, không thể cãi được.''

Bỗng cchoosc, một luồng sáng lóe lên ngay giữa khoang xe. Không cánh cửa, không báo trước, ``vị thần'' Choco - hay ai đó giống hệt cô - xuất hiện từ hư vô, đứng lơ lửng giữa hai hàng ghế. Cô khoanh tay, nở nụ cười bí ẩn, như thể đã theo dõi toàn bộ sự việc từ đầu.

``Có vẻ như hai con cần chút giúp đỡ nhỉ.''

\img{5-3k}

Trước sự xuất hiện bất ngờ của cô, nàng Sói đen hoảng hồn, lắp bắp: ``V-Vị thần lúc nãy!? Sao chị lên xe được? Xe đang chạy với tốc độ Mach 100 cơ mà!''

Kuon thì chẳng còn bất ngờ nữa. Cậu chỉ thở dài, giọng mệt mỏi như đã đoán định được tình huống: ``Chắc lại toán nữa chứ gì? Lại phương trình, lại ẩn số, lại thời gian đuổi kịp?''

Choco gật đầu đầy thần bí, sau đó đột nhiên giơ một tay lên, tạo hình chữ C bằng hai ngón tay, và dõng dạc tuyên bố, giọng vang vọng khắp xe: ``Thực ra, ta là Tiên XiP3 - gọi là điểm Xi\footnotemark - di chuyển với vận tốc Mach 100.''

\img{5-3l}
\footnotetext{Hai cụm từ vần với nhau, lần lượt mang nghĩa là ``thiên sứ'' (天使) và ``điểm C'' (点C).}

Ngay lập tức, Korone và Aki ở phía sau - vốn dĩ chẳng hiểu chuyện gì đang xảy ra - cũng bắt chước làm theo. Họ giơ tay tạo thành chữ C, nháy mắt tạo dáng, nở một nụ cười thật tươi, như thể đang tham gia một nghi lễ bí truyền.

``Tiên XiP3!'' Korone reo lên, miệng vẫn nhai kẹo.

``Tiên XiP3!'' Aki cười toe, giơ điện thoại selfie lên chụp luôn khoảnh khắc này.

Subaru ngồi giữa, chỉ biết há hốc mồm nhìn họ như thể vừa bước vào một giáo phái bí ẩn. Cô lắp bắp: ``C-Các cậu bị gì thế? Có ai giải thích cho mình hiểu không? Sao tự dưng lại biến thành\dots\ tiên?''

Mio sững sờ, gương mặt méo mó vì không hiểu nổi chuyện gì đang xảy ra. Cô đưa tay ôm đầu, than: ``Thế nghĩa là sao chứ!? Tiên XiP3 là cái quái gì? Và sao ai cũng làm theo hết vậy!? Đây có phải là tà giáo không?''

Cậu con trai đưa tay xoa trán, giọng bất lực tột độ: ``Anh cũng chịu em ạ\dots\ Chắc là do không khí trong xe loãng quá, thiếu oxy nên mọi người bị ảo giác tập thể rồi.''

Rồi, cậu quay sang Okayu, người vẫn đang thư thái lái xe, giương ánh mắt tuyệt vọng: ``Ai đó cho anh xuống xe cái\dots\ Chứ ngồi đây lâu nữa chắc anh cũng biến thành Tiên XiP3 luôn quá.''

%==========================================TRUYỆN PHỤ=========================================
\omk

\pov{Hikawa Kuon}

Thực ra, bài toán của Choco bị thiếu dữ kiện. Nguyên văn bài toán đó như sau:

\begin{tcolorbox}[
    colframe=vol,    % Màu viền
    colback=white,   % Màu nền
    boxrule=0.5pt,   % Độ dày viền
    arc=0mm,         % Góc vuông (không bo tròn)
    left=5pt,        % Khoảng cách giữa nội dung và viền trái
    right=5pt,       % Khoảng cách giữa nội dung và viền phải
    top=5pt,         % Khoảng cách giữa nội dung và viền trên
    bottom=5pt       % Khoảng cách giữa nội dung và viền dưới
  ]
  Matsuri (gọi là A) rời khỏi văn phòng, đi với tốc độ Mach 80 được 10 phút thì dừng lại. Tôi và Mio (gọi là B) đuổi theo, nhưng chỉ chạy được 2km với tốc độ 10km/h thì được Okayu-chan được lên xe đi nhờ, chạy với vận tốc Mach 100. Tính thời gian mà B bắt kịp được A. Gọi Mach 1 tương đương với 340m/s.
\end{tcolorbox}

Cái điểm C (đọc là điểm Xi) đó thực chất là chơi chữ, nên chẳng liên quan gì đến bài toán cả.

Các bạn thử tính xem chúng tôi đã mất bao lâu để bắt kịp cô nàng rơi não này. Đáp án tôi sẽ để ở bên dưới phần {\abv Đáp án}.

\il

Sau một hồi tính toán loạn xạ, những phương trình nảy số trong đầu, và những cú phóng như bay trên chiếc xe không biết gọi là gì cho chính xác - tạm gọi là ``tên lửa Lamb*rgh*ni'' cuối cùng chúng tôi cũng bắt kịp Matsuri.

Okayu giảm tốc độ, chiếc xe từ từ hạ xuống, lơ lửng cách mặt đất chưa đầy nửa mét. Mio và tôi nhanh chóng nhảy ra ngoài, chạy lại chỗ nàng Cổ động đang đứng như một pho tượng giữa đường, mắt vẫn còn đờ đẫn sau cú ``tên lửa vụt qua'' lúc nãy.

``Của chị đây! Điện thoại, ví, thẻ, chìa khóa - tất cả!'' Mio vừa thở hổn hển vừa dúi đống đồ vào tay Matsuri, giọng như sắp khóc vì sung sướng.

Matsuri chớp mắt nhìn đống đồ, rồi ngước lên nhìn chúng tôi. Ngay khi nhận ra mình vừa được ``cứu trợ khẩn cấp'', nàng Lễ hội như được hồi sinh. Khuôn mặt cô sáng bừng, mắt cười tít lại, miệng nở nụ cười rạng rỡ như ánh mặt trời: ``Wa! Cảm ơn hai người nhiều nha! Tưởng hôm nay phải ngủ ngoài đường rồi!''

Rồi cô ôm chầm lấy Mio, suýt làm cô ngã. Tôi đứng bên cạnh, cười nhạt, thở phào nhẹ nhõm. Cuộc đuổi bắt xem như đã kết thúc có hậu.

Thế nhưng\dots\ Ngay khi đang vui vẻ, Matsuri bỗng dưng ngừng lại, nhìn đống đồ trên tay, rồi nhìn chúng tôi, rồi lại nhìn đống đồ. Cô ngơ ngác nói một câu khiến cả tôi và Mio chết lặng: ``Ủa\dots\ em không cần có nhiều đồ đến thế đâu.''

``Hả? Là sao?'' cả hai chúng tôi đồng thanh, mặt méo xệch.

Matsuri tiếp tục, giọng đầy vẻ ngây thơ đến tội nghiệp: ``Em chỉ ra ngoài mua chút đồ ăn thôi mà\dots\ Ở tiện lợi cách công ty có năm trăm mét. Em chỉ quên mỗi ví với điện thoại ở văn phòng thôi. Chìa khóa nhà em vẫn mang theo đây nè.''

Cô rút từ túi quần ra một chiếc chìa khóa, kẹp cùng một chiếc móc hình cà rốt nhỏ xíu.

``Thế còn đống đồ trên bàn?'' Mio hỏi, giọng đã bắt đầu run.

``À, mấy cái đó là đồ để trong balo em từ hôm qua, tối về em mới cất. Chứ có cần mang đi đâu.''

Nghe câu trả lời của em ấy, mặt Mio đỏ bừng - không phải vì trời nắng gắt, mà vì tức. Cô quay lưng đi, đập tay lên thành xe của Okayu một cái ``bộp'', vừa lẩm bẩm chửi đổng vừa tức đến rơi nước mắt: ``Thế mà Fubuki lại bảo tôi là quên nhiều đồ lắm\dots\ `Mang hết đồ cho nó đi, nó sắp chết đói giữa đường rồi!'\dots\ `Nó đang đứng như cá ươm ở đâu đó, mau đi cứu!'\dots\ Thật là\dots\ Mất công bọn tôi phải mang cả đống đồ lỉnh kỉnh chạy bộ, rồi ngồi tên lửa, rồi giải toán Mach, rồi phi như bay\dots\ để đưa cho cô một đống đồ cô chẳng cần!''

Matsuri ngơ ngác nhìn Mio, rồi quay sang tôi, như muốn hỏi ``có chuyện gì vậy?''.

Tôi chỉ có thể cười trừ, vỗ vai Mio để an ủi: ``Thôi, không sao. Đưa được đồ cho em ấy là may rồi. Ít nhất thì em ấy đã có điện thoại và ví, mai khỏi phải quay lại văn phòng để lấy.''

Mio vẫn sụt sịt, mặt cúi gằm. Nàng Sói đen tội nghiệp đã trải qua một buổi sáng quá sức chịu đựng.

\ct

Sau khi Matsuri - người mà từ nay tôi sẽ gọi thầm là ``não cá vàng'' mua đồ xong (một túi bánh snack và hai chai nước suối), tất cả chúng tôi kéo nhau trở lại xe để về văn phòng.

Tôi tưởng rằng mọi chuyện sẽ yên ả, êm đềm. Nhưng tôi đã quên mất một sự thật hiển nhiên của vũ trụ: khi có từ ba cô gái trở lên trong một không gian chật hẹp, bình yên không tồn tại.

Bầu không khí trên xe ngay sau đó trở nên vô cùng hỗn loạn. Trong khi tôi và Mio vẫn còn ngồi im với vẻ mặt bất lực, như hai xác ướp vừa được khai quật, thì những thành viên còn lại - ai nấy đều rất vui vẻ. Quá vui là đằng khác.

\img{5-3m}

Korone thì không biết bằng cách nào đã trèo từ hàng ghế giữa sang chỗ Okayu và Mio. Chưa dừng lại ở đó, em ấy còn leo lên hẳn ``cân đẩu vân'' của Choco-san - tức là cái đám mây nhỏ mà ``vị thần'' đang đứng. Trông cô nàng lúc này chẳng khác gì một phi hành gia đang thực hiện cuộc đi bộ ngoài không gian.

``Ogayu, cho tôi lái thử với\~{}'' Korone nắm tay Okayu, mắt long lanh cầu xin, giọng ngọt như mía lùi.

Okayu - nàng Mèo điềm nhiên - vẫn ung dung giữ vô lăng, đáp lại bằng giọng đều đều: ``Lần sau nhé bạn hiền\~{} Bây giờ đang ở tốc độ Mach 80, lỡ tay cái là bay lên trời luôn đấy.''

Ngay cạnh đó, Mio đổ mồ hôi hột nhìn sự vô tư của nàng Chó. Cô lên tiếng nhắc nhở, giọng đầy lo lắng: ``Korone, đừng có sà xuống bọn này chứ! Nguy hiểm lắm đó! Với lại trả lại đồ cho Choco-sensei đi! Cậu đang đứng trên mây của người ta kìa!''

Korone giả vờ rầu rĩ, bĩu môi, đáp lại bằng một giọng tinh nghịch: ``Ơ kìa\~{} Đang vui mà\~{} Mây của Choco-sensei êm lắm, không bị rơi đâu.''

Trong khi đó, ở hàng ghế giữa, Matsuri và Subaru đang\dots\ giật đồ nghề của nàng Y tá. Chính xác, họ đã cuỗm mất ``bảo bối'' của ``vị thần'' ngay trước mặt cô.

Matsuri cầm lấy chiếc băng-rôn màu trắng có chữ ``Thần'' mà Choco đeo trên đầu (cái mà cô dùng để tạo điểm nhấn cho vai diễn ``Tiên XiP3''), quấn lên trán mình, tỏ vẻ thích thú: ``Cái băng-rôn này trông cũng hay đấy! Cho em đeo thử với! `Thần' này rất hợp với em luôn!''

Subaru thì cuỗm mất chiếc gậy trượng của ``vị thần''. Cô nàng Vịt cầm gậy lên, đôi mắt sáng rực như thể đây là lần đầu tiên trong đời em ấy nhìn thấy một cây gậy ma thuật: ``Em muốn dùng thử gậy phép này! `Biến hình!' \dots Mà Choco-sensei, làm thế nào để nó phun ra pháo hoa vậy ạ?''

Choco - lúc này đã bị tước mất cả mây, cả băng, cả gậy - chỉ còn biết đứng lơ lửng bằng\dots\ uy lực thần thánh (hoặc bằng đôi chân trần). Cô vừa cố giành lại đồ vừa kêu lên, giọng đầy bất lực: ``Không được đâu, Matsuri-sama, Subaru-sama! Trả lại cho ta đi! Cái băng-rôn là bảo vật gia truyền, cây gậy là pháp khí thiêng liêng! Trả lại đây!''

Nhưng cả hai cô gái đều cười khúc khích, chạy lên chạy xuống hàng ghế, biến chiếc xe tên lửa chật hẹp thành một sân khấu hài kịch di động.

Aki đứng nép vào một góc, tay cầm điện thoại selfie, vừa quay vừa cười, trông như một khán giả đang thưởng thức một vở kịch miễn phí: ``Trông mọi người có vẻ vui nhỉ. Mà Choco-sensei đóng vai làm thần ạ? Em thấy cô đóng hay lắm, nhất là cái đoạn `Ta là Tiên XiP3' ấy.''

Tôi ngồi hàng ghế cuối, tặc lưỡi, thở dài bất lực, đưa tay xoa trán: ``Thôi, hết cứu. Mới sáng ra đã chạy đuổi người bằng tên lửa, xong giờ lại biến xe thành rạp xiếc. Mệt với mấy người quá.''

``Cái câu đó phải là em nói mới đúng đó, Kuon-senpai\dots'' Mio thẫn thờ đáp lại, ánh mắt còn chẳng thèm ngoái nhìn lại đằng sau tôi.

Chiếc xe ``Lamb*rgh*ni tên lửa'' lướt nhẹ trên bầu trời, chở theo một đám người hỗn độn nhưng đầy tiếng cười. Bên ngoài cửa sổ, mây trôi, gió thổi, nắng vàng trải dài. Matsuri đang nghịch băng-rôn ``Thần'', Subaru huơ huơ cây gậy, Korone vẫn đeo bám Okayu để xin lái xe, Mio thì đã nằm gục ra ghế, mệt mỏi đến nỗi chẳng buồn cãi nhau nữa. Còn Choco? Cô đã giành lại được cây gậy và băng-rôn sau một hồi đuổi bắt, nhưng chiếc mây thì vẫn bị Korone chiếm giữ.

Aki vẫn quay video, cười khoái chí. Tôi nhìn tất cả, lắc đầu, khẽ cười.

``Hầy\dots\ Một ngày của cả một bầu trời vận tốc là như vậy đấy, mọi người ạ. Vừa chạy bộ, vừa giải toán, vừa ngồi tên lửa, vừa xem hài kịch. Chẳng biết mai còn gì nữa\dots''

\ct

Xe từ từ hạ cánh xuống bãi đỗ trước văn phòng. Cửa mở ra, từng người từng người bước ra, mặt mũi hớn hở.

Mio cuối cùng cũng thở được một hơi thật sâu: ``Em về đây\dots\ Em xin thề, lần sau có ai quên đồ, người đó tự lo. Em không đi giao hàng tốc hành nữa.''

Matsuri chạy lại, khoác vai Mio, cười tươi: ``Cảm ơn Mio-chan nhé! Lát em mua trà sữa mời!''

Mio chỉ lắc đầu, bất lực mỉm cười.

Tôi đi sau cùng, nhìn bóng lưng họ. Một buổi sáng đầy mệt mỏi nhưng cũng tràn ngập những khoảnh khắc đáng nhớ.

\begin{center}{\abv\Large Đáp án bài toán}\end{center}

Trước tiên, ta phải quy đổi đơn vị phép tính về km/h:
\begin{itemize}
  \item Mach 1 = 340m/s = 1224km/h
  \item Mach 80 = \(1224 \times 80 = 97920\) km/h
  \item Mach 100 = \(1224 \times 100 = 122400\)km/h
\end{itemize}

Sau đó, ta sẽ tính quãng đường mà Mat-chan đã đi. Do em ấy đi với tốc độ Mach 80 trong 10 phút (tức là \(\frac{1}{6}\) giờ) nên:
\begin{align*}
  S_A = v_A \times t_A = 97920 \times \frac{1}{6} = 16320 (km)
\end{align*}
Sau đó, em ấy dừng lại, tức là sẽ đứng yên tại điểm này.

Tiếp theo, ngay khi em ấy dừng, tôi và Mio-chan (có thể coi là) bắt đầu đuổi theo với quãng đường 2km, vận tốc 10km/h, tức là:
\begin{align*}
  t_B = \frac{2}{10} = 0.2 (h)
\end{align*}
Sau khi chúng tôi chạy được trong 2 giờ (= 12 phút), chúng tôi được Okayu-chan chở đi. Thời điểm đó, chúng tôi chỉ còn lại:
\begin{align*}
  D = 16320 - 2 = 16318 (km)
\end{align*}
Chiếc xe của Okayu-chan đạt vận tốc Mach 100, nên thời gian chúng tôi bắt kịp được em ấy sẽ là:
\begin{align*}
  t = \frac{16318}{122400} \approx 0.13333 (h)
\end{align*}
hay 8 phút.

Vậy tổng thời gian chúng tôi chạy để bắt kịp Matsuri-chan sẽ là:
\begin{align*}
  T = 12 + 8 = 20 (m)
\end{align*}
Như vậy, sau 20 phút chúng tôi đã đuổi kịp được em ấy.

\end{document}